\figref{eval-example} shows three rules of the expression judgement, which stand for the rest.
\ruleName{eval-binop} evaluates both operands and binds the outcome of the second. The meaning
$\hat{\odot}(v, v')$ of the operator is itself an outcome, $\aborts{\kappa}$ where Python raises, and is
undefined where \PurePy gives the operation no meaning; an expression whose operator meaning is undefined
has no derivation. \ruleName{eval-binop-fail} propagates an abort from the first operand.
\ruleName{eval-call-nonfun} originates an abort: calling a value that is neither closure nor primitive
is a \pyinline{TypeError}, as in Python. The remaining rules follow the same pattern and are given in
full in \specification.

\begin{figure}
\begin{mathpar}
\inferrule*[lab={\ruleName{eval-binop}}]
{
  \rho, e \Rightarrow \val{v} \\
  \rho, e' \Rightarrow r
}
{
  \rho, \exBinOp{e}{\odot}{e'} \Rightarrow \mlet{v'}{r}{\hat{\odot}(v, v')}
}

\and
\inferrule*[lab={\ruleName{eval-binop-fail}}]
{
  \rho, e \Rightarrow \aborts{\kappa}
}
{
  \rho, \exBinOp{e}{\odot}{e'} \Rightarrow \aborts{\kappa}
}

\and
\inferrule*[lab={\ruleName{eval-call-nonfun}}]
{
  \rho, e \Rightarrow \val{v} \\
  \rho, \seq{e} \Rightarrow \val{\seq{u}} \\
  v \text{ neither closure nor primitive}
}
{
  \rho, \exCall{e}{\seq{e}} \Rightarrow \aborts{\errType}
}

\end{mathpar}
\caption{Expression evaluation: representative rules}
\label{fig:eval-example}
\end{figure}
